\documentclass[a4paper,landscape, 12pt]{extarticle}
\usepackage[left=0.5in, right=0.5in, top=0.4in, bottom=0.55in]{geometry}
\usepackage{graphicx} % Required for inserting images
\usepackage[normalem]{ulem}
\usepackage{tabto}
\usepackage{minted}
\usepackage{float}
\renewcommand\ULthickness{1.0pt}   %%---> For changing thickness of underline
\setlength\ULdepth{1.3ex}%\maxdimen ---> For changing depth of underline
\usepackage{tikz}
\usetikzlibrary{graphs}

\usepackage{array}
\usepackage{caption}
\usepackage{amsmath} 
\usepackage{fancyhdr}
\usepackage{lastpage}
\usepackage{enumitem}

\author{Juan R Ortiz; Timothy Jung}

% Page style with slide numbers
\pagestyle{fancy}
\fancyhf{}
\renewcommand{\headrulewidth}{0pt}
\fancyfoot[R]{\Large\thepage\ / \pageref{LastPage}}

% Define header command for consistent formatting
\newcommand{\slideheader}[2]{%
\noindent

\includegraphics[height=0.5in]{../../images/cpp_logo.png}%
\hspace{0.3in}%
{\Large \uline{#1 \hfill #2}}%
\par\vspace{0.4in}
}
% Standard content spacing - 0.6in base indent for all content
\newcommand{\contentstart}{%
\leftskip 0.6in \rightskip 0.6in
}
\newcommand{\contentend}{%
\leftskip 0in \rightskip 0in
}

% Standard font size for body content
\newcommand{\bodytext}{\LARGE}

\begin{document}
\begin{huge} 
%%%%%%%%%%%%%%%%%%%%%%%%%%%%%%%%%%%%%%%%
% SLIDE 1 - Title Slide
\slideheader{ECE5110 Numerical Modeling}{Unit 05}

\vspace{1.2in}
\begin{center}
{\Huge \textbf{Juan Ortiz and Timothy Jung}}

\vspace{0.4in}
{\LARGE Modeling and Simulation of a Temperature Sensor Using the Runge–Kutta Method}
\end{center}
%%%%%%%%%%%%%%%%%%%%%%%%%%%%%%%%%%%%%%%%
\newpage %Page 2
\slideheader{Unit 05: Ordinary Differential Equations}{ECE5110 - Numerical Modeling}

\leftskip 0.7in \rightskip 0.7in 
\vspace{0.4in}
\noindent Despite temperature sensors being in controlled environments, they can still be affected by external noise, such electromagnetic interference or mechanical vibrations. 
\vspace{0.1in}

We will model a noisy temperature sensor's output response to a step change in temperature using an ordinary differential equation (ODE).
\vspace{0.1in}

\noindent First order ODE form:

\[\hspace{0.4in}
\frac{dT}{dt} = -a(T - T_{ambient}) + K_{p}(T_{setpoint} - T)
\]

\begin{center}
\begin{minipage}{0.75\linewidth} % adjust width as you like
\begin{minted}[fontsize=\Large, frame=single]{python}
def f(t, y):
    return -a * (y - T_ambient) 
                  + Kp * (T_set - y)
\end{minted}
\end{minipage}
\end{center}


\vspace{0.2in}

\leftskip 0in \rightskip 0in 

%%%%%%%%%%%%%%%%%%%%%%%%%%%%%%%%%%%%%%%%
\newpage %Page 3
\slideheader{Unit 05: Ordinary Differential Equations}{ECE5110 - Numerical Modeling}

\leftskip 0.7in \rightskip 0.7in 
\vspace{0.4in}
\noindent Each term in the ODE has a specific physical meaning:
\vspace{0.4in}

\begin{minipage}[t]{\dimexpr\textwidth-1.2in}
\begin{itemize}
    \item $T$: sensor temperature
    \item $t$: time
    \item $a$: positive cooling rate constant
    \item $T_{ambient}$: ambient temperature
    \item $K_{p}$: proportional gain constant
    \item $T_{setpoint}$: desired temperature
\end{itemize}
\end{minipage}


\vspace{0.2in}

\leftskip 0in \rightskip 0in 

%%%%%%%%%%%%%%%%%%%%%%%%%%%%%%%%%%%%%%%%
\newpage %Page 4
\slideheader{Unit 05: Ordinary Differential Equations}{ECE5110 - Numerical Modeling}

\leftskip 0.7in \rightskip 0.7in 
\vspace{0.4in}
\noindent We can solve this ODE by implementing the 4th order Runge-Kutta method.
\vspace{0.4in}

\noindent RK4 advances the solution using: 
\begin{align*}
k_1 &= f(t_n, y_n) \\
k_2 &= f(t_n + \frac{h}{2}, y_n + \frac{h}{2}k_1) \\
k_3 &= f(t_n + \frac{h}{2}, y_n + \frac{h}{2}k_2) \\
k_4 &= f(t_n + h, y_n + hk_3) \\
y_{n+1} &= y_n + \frac{h}{6}(k_1 + 2k_2 + 2k_3 + k_4)
\end{align*}




\leftskip 0in \rightskip 0in 

%%%%%%%%%%%%%%%%%%%%%%%%%%%%%%%%%%%%%%%%

\newpage %Page 5
\slideheader{Unit 05: Ordinary Differential Equations}{ECE5110 - Numerical Modeling}

\leftskip 0.7in \rightskip 0.7in 
\vspace{0.4in}

\noindent This can be implemented in Python as follows:: 
\begin{center}
\begin{minipage}{0.75\linewidth} % adjust width as you like
\begin{minted}[fontsize=\huge, frame=single]{python}
for i in range(1, len(steps)):
  x_prev = steps[i - 1]
  y_prev = y[i - 1]
  k1 = f(x_prev, y_prev)
  k2 = f(x_prev + h / 2, 
         y_prev + h * k1 / 2)
  k3 = f(x_prev + h / 2, 
         y_prev + h * k2 / 2)
  k4 = f(x_prev + h, 
         y_prev + h * k3)
  y[i] = y_prev + h * (k1 + 2 * k2 + 
                            2 * k3 + 
                                k4) / 6
\end{minted}
\end{minipage}
\end{center}

\leftskip 0in \rightskip 0in 

%%%%%%%%%%%%%%%%%%%%%%%%%%%%%%%%%%%%%%%%

\newpage %Page 5
\slideheader{Unit 05: Ordinary Differential Equations}{ECE5110 - Numerical Modeling}

\leftskip 0.7in \rightskip 0.7in 
\vspace{0.4in}

\noindent To simulate the sensor noise, we add a random noise term to the ODE output at each time step:

\begin{center}
\begin{minipage}{0.75\linewidth} % adjust width as you like
\begin{minted}[fontsize=\huge, frame=single]{python}
noise = np.random.normal(0.0, 
                         noise_std, 
                         size=len(y)
                         )
y_noisy = y + noise
\end{minted}
\end{minipage}
\end{center}

\leftskip 0in \rightskip 0in 

%%%%%%%%%%%%%%%%%%%%%%%%%%%%%%%%%%%%%%%%
\newpage %Page 6
\slideheader{Unit 05: Ordinary Differential Equations}{ECE5110 - Numerical Modeling}

\leftskip 0.7in \rightskip 0.7in 
\vspace{0.4in}

\noindent The plot below shows the simulated temperature response of the sensor to a step change in temperature:

\begin{center}
\includegraphics[height=5.0in]{../../images/unit05_rk4_noisy.png}%
\end{center}

\leftskip 0in \rightskip 0in 

%%%%%%%%%%%%%%%%%%%%%%%%%%%%%%%%%%%%%%%%
\newpage %Page 7
\slideheader{Unit 05: Ordinary Differential Equations}{ECE5110 - Numerical Modeling}

\leftskip 0.7in \rightskip 0.7in 
\vspace{0.4in}

\noindent This model is practical in industry because it approximates real sensor behavior.
\vspace{0.4in}

\begin{minipage}[t]{\dimexpr\textwidth-1.2in}
\begin{itemize}
    \item It allows engineers to predict temperature response.
    \item Helps tune control parameters.
    \item Supports testing before hardware deployment.
    \item Reduces cost and development time.
    \item Improves system reliability.

\end{itemize}
\end{minipage}

\leftskip 0in \rightskip 0in 

%%%%%%%%%%%%%%%%%%%%%%%%%%%%%%%%%%%%%%%%

\newpage %End Page
\slideheader{Unit 05: Ordinary Differential Equations}{ECE5110 - Numerical Modeling}

\leftskip 0.7in \rightskip 0.7in 
\vspace{2.5in}
\begin{center}
{\Huge \textbf{QUESTIONS?}}
\end{center}
\leftskip 0in \rightskip 0in 


%%%%%%%%%%%%%%%%%%%%%%%%%%%%%%%%%%%%%%%%


\end{huge}
\end{document}